\documentclass[../main.tex]{subfiles}
\begin{document}
\subsection{F1 and F2: execution-cost atlas and engine validation}\label{sec:f1f2}

\paragraph{Economic intuition and literature.}
Quoted spreads can overstate the costs that liquidity takers pay \citep{muravyevpearson2020}. Wave 2 therefore began with a measurement study (F1) and an engine check (F2). Neither is an alpha test. If the engine cannot reproduce a published, mechanically defined index, its strategy P\&L is not credible.

\paragraph{Specification.}
\begin{itemize}
\item \textbf{F1 (measurement).} For each trade with a valid NBBO, compute the effective half-spread ratio and the relative quoted spread,
\begin{equation}
\rho = \frac{|p - m|}{(a - b)/2}, \qquad s_{\mathrm{rel}} = \frac{a - b}{m},
\end{equation}
where $p$ is the trade price, $m$ the NBBO mid, $a$ the ask and $b$ the bid. Results are bucketed by root, DTE $\{0, 1\text{--}7, 8\text{--}30, 31\text{--}90, >90\}$, moneyness $|K/S - 1| \in \{<1\%, 1\text{--}3\%, 3\text{--}7\%, >7\%\}$ and half-hour, with $S$ the prior close. A timing check compares the EOD-chain bid/offer with the tape's last NBBO before 16:00 and 16:15 for SPY, SPXW and AAPL.
\item \textbf{F2 (replication).} On each third-Friday expiry, sell one SPXW PM put expiring next month. The strike is the listed strike nearest to and not above the SPX close. The fill is the EOD bid (T3; the mid is a sensitivity). Collateral is the full strike notional in T-bills. Hold to expiry and settle on the SPX close (T0).
\item \textbf{Execution tier:} T3 entry and T0 exit for F2. \textbf{Risk:} the put is fully cash-secured.
\end{itemize}

\paragraph{Implementation.}
F1 uses \file{wave2/extract.py} to build the atlas from about 1.5 billion sampled development-period prints. F2 is run by the common engine (\file{wave2/cpp/common.hpp} and the F2 engine). The F2 tape sensitivity uses the first bid within 120 s of the 15:55 sale (amendment A3). The primary EOD variant is unchanged. Differences from the Cboe PUT index are expected because PUT uses AM settlement and a VWAP sale window. These differences are documented rather than fitted.

\paragraph{Out-of-sample testing and results.}
Neither study has a holdout test. For F1, SPXW prints well inside the spread, SPY and QQQ print near the touch, and single names fall in between (\figref{fig:cost}{fig02_cost_atlas.png}). The effective/quoted half-spread ratio is 0.57 at 0DTE ATM and 0.34 at 8--30 DTE for SPXW, and 0.75--0.93 for SPY and QQQ. Other traders' effective spreads are not our fills, so T2 still crosses the full spread.

For F2, the monthly cash-secured SPXW put-write tracks the Cboe PUT index (Table~\ref{tab:f2}). The mean gap is in line with PUT's AM settlement and VWAP sale window, which we do not replicate. We read this as evidence that the option accounting is sound.

\begin{table}[tbp]
\centering\small
\caption{F2: put-write engine against the Cboe PUT index.}\label{tab:f2}
\begin{tabular}{@{}lr@{}}
\toprule
Statistic & Value \\
\midrule
Monthly return correlation & 0.972 \\
Tracking error & 2.3\% per year \\
Mean monthly gap & $-$9 bp \\
Monthly rolls & 34 \\
\bottomrule
\end{tabular}
\end{table}

\Fig{wave2}{fig03_putwrite_validation.png}{F2: put-write engine against the Cboe PUT index (\texttt{wave2/figures/fig03\_putwrite\_validation.png}).}{fig:f2}
\ifSubfilesClassLoaded{\section*{References}% Bibliography.
\bigskip\noindent\emph{The reference details below should be checked against the published versions before citing.}\par\medskip
\begin{thebibliography}{99}
\setlength{\itemsep}{2pt}
\small
\bibitem[Bailey and L{\'o}pez de Prado(2014)]{bailey2014} Bailey, D. H., \& L{\'o}pez de Prado, M. (2014). The deflated Sharpe ratio. \emph{Journal of Portfolio Management}, 40(5).
\bibitem[Barber et~al.(2013)]{barber2013} Barber, B. M., De George, E. T., Lehavy, R., \& Trueman, B. (2013). The earnings announcement premium around the globe. \emph{Journal of Financial Economics}, 108(1).
\bibitem[Brandt et~al.(2008)]{brandt2008} Brandt, M. W., Kishore, R., Santa-Clara, P., \& Venkatachalam, M. (2008). Earnings announcements are full of surprises. Working paper.
\bibitem[Chan et~al.(1996)]{chan1996} Chan, L. K. C., Jegadeesh, N., \& Lakonishok, J. (1996). Momentum strategies. \emph{Journal of Finance}, 51(5).
\bibitem[Cremers and Weinbaum(2010)]{cremers2010} Cremers, M., \& Weinbaum, D. (2010). Deviations from put-call parity and stock return predictability. \emph{Journal of Financial and Quantitative Analysis}, 45(2).
\bibitem[DeMiguel et~al.(2009)]{demiguel2009} DeMiguel, V., Garlappi, L., \& Uppal, R. (2009). Optimal versus naive diversification: How inefficient is the 1/N portfolio strategy? \emph{Review of Financial Studies}, 22(5).
\bibitem[Easley et~al.(1998)]{easley1998} Easley, D., O'Hara, M., \& Srinivas, P. S. (1998). Option volume and stock prices: Evidence on where informed traders trade. \emph{Journal of Finance}, 53(2).
\bibitem[Fama and MacBeth(1973)]{famamacbeth1973} Fama, E. F., \& MacBeth, J. D. (1973). Risk, return, and equilibrium: Empirical tests. \emph{Journal of Political Economy}, 81(3).
\bibitem[Frazzini and Lamont(2007)]{frazzini2007} Frazzini, A., \& Lamont, O. A. (2007). The earnings announcement premium and trading volume. NBER Working Paper 13090.
\bibitem[Johnson and So(2012)]{johnsonso2012} Johnson, T. L., \& So, E. C. (2012). The option to stock volume ratio and future returns. \emph{Journal of Financial Economics}, 106(2).
\bibitem[Jorion(1986)]{jorion1986} Jorion, P. (1986). Bayes-Stein estimation for portfolio analysis. \emph{Journal of Financial and Quantitative Analysis}, 21(3).
\bibitem[Ledoit and Wolf(2004)]{ledoit2004} Ledoit, O., \& Wolf, M. (2004). A well-conditioned estimator for large-dimensional covariance matrices. \emph{Journal of Multivariate Analysis}, 88(2).
\bibitem[Markowitz(1952)]{markowitz1952} Markowitz, H. (1952). Portfolio selection. \emph{Journal of Finance}, 7(1).
\bibitem[Newey and West(1987)]{neweywest1987} Newey, W. K., \& West, K. D. (1987). A simple, positive semi-definite, heteroskedasticity and autocorrelation consistent covariance matrix. \emph{Econometrica}, 55(3).
\bibitem[Pan and Poteshman(2006)]{pan2006} Pan, J., \& Poteshman, A. M. (2006). The information in option volume for future stock prices. \emph{Review of Financial Studies}, 19(3).
\bibitem[Xing et~al.(2010)]{xing2010} Xing, Y., Zhang, X., \& Zhao, R. (2010). What does the individual option volatility smirk tell us about future equity returns? \emph{Journal of Financial and Quantitative Analysis}, 45(3).
\end{thebibliography}
}{}
\end{document}
